Die Gliederung der Arbeit erfolgt mit den gewöhnlichen
\LaTeX-Gliederungsbefehlen \verb+\chapter+, \verb+\section+,
\verb+\subsection+, usw. Vergleiche hierzu auch die
Listings~\ref{lst:tpl:article} bis~\ref{lst:tpl:book}. Die \texttt{*}-Varianten
werden nicht nummeriert und erscheinen standardmäßig nicht im
Inhaltsverzeichnis. Möglichkeiten zur Anpassung des Inhaltsverzeichnisses
werden in Abschnitt~\ref{sec:final:toclists} behandelt.

Für die Gliederung ist folgendes Musterschema zweckmäßig. Je nach Art der
Arbeit entfallen die nicht zutreffenden Punkte.

\begin{compactitem}
\item Deckblatt
\item Widmung
\item Kurzreferat
\item Aufgabenstellung
\item Inhaltsverzeichnis
\item weitere Verzeichnisse
      (ggf. auf am Ende der Arbeit möglich, siehe unten)
\item Abkürzungsverzeichnis, Symbole und Bezeichner
\item \textbf{Textkörper}
\item Anhänge
\item Literaturverzeichnis
\item weitere Verzeichnisse
      (Alternative, sofern nicht bereits am Anfang gesetzt)
\item Glossar
\item Index
\item Danksagung
\item Selbständigkeitserklärung
\end{compactitem}
